% ============================================================================
% 05_ers_ieee830.tex — Cuerpo del entregable "ERS IEEE Std 830-1998"
% ============================================================================
\section{Introducción}

\subsection{Propósito}

El propósito de este documento es especificar, de manera completa,
verificable y trazable, los requisitos funcionales y no funcionales del
sistema \textbf{LSA Salud}, un prototipo de software para la comunicación
bidireccional entre pacientes que usan Lengua de Señas Argentina (LSA) y
personal de recepción o enfermería, pensado para el Sanatorio 9 de Julio,
en entornos de atención médica donde no hay un intérprete humano
disponible. Sigue la estructura recomendada por el estándar IEEE Std
830-1998. Los requisitos oficiales (RF-01 a RF-06, RNF-01 a RNF-07) provienen
de la matriz de trazabilidad de requisitos de la cátedra (TP2); este
documento los desarrolla en sub-requisitos técnicos trazables (ver
Sección~3.2 y \texttt{IDS.md}).

Este documento está dirigido a:
\begin{itemize}
    \item El equipo de desarrollo del proyecto.
    \item Docentes y evaluadores académicos de la UTN--FRT.
    \item Profesionales de salud interesados en evaluar el prototipo.
    \item Futuros colaboradores o integradores del sistema.
\end{itemize}

\subsection{Ámbito del sistema}

\textbf{Nombre del sistema:} LSA Salud --- Sistema Inteligente de
Comunicación y Reconocimiento de LSA. \textbf{Cliente:} Sanatorio 9 de
Julio (San Miguel de Tucumán). El nombre del repositorio de código,
``Nexora'', es un alias técnico interno y no el nombre del proyecto
académico.

\subsubsection{Lo que el sistema hace}
\begin{itemize}
    \item Captura en tiempo real los movimientos corporales del paciente
    mediante visión por computadora (MediaPipe Holistic o MediaPipe
    Tasks).
    \item Clasifica las señas capturadas con una red neuronal LSTM
    entrenada con datos propios.
    \item Reconstruye oraciones naturales en español a partir de las
    palabras sueltas reconocidas, usando plantillas médicas precalculadas
    y, como respaldo, un modelo de lenguaje local.
    \item Sintetiza las oraciones en voz con acento argentino.
    \item Provee un canal de comunicación bidireccional entre el paciente
    y el personal de recepción o enfermería, sincronizado entre dos estaciones de trabajo.
    \item Permite al personal de recepción o enfermería responder por
    texto o por dictado de voz.
    \item Ofrece un módulo de captura de datos para ampliar el vocabulario
    del modelo y un script de reentrenamiento.
    \item \textit{Diseñado, pendiente de implementación:} solicita
    consentimiento explícito de cámara antes de iniciar la captura
    (RF-03.1), registra eventos de auditoría anónimos por operación
    relevante (RF-04.1) y permite reportar una traducción incorrecta o
    ambigua (RF-06.1). Ver Sección~3.2 y \texttt{IDS.md}.
\end{itemize}

\subsubsection{Lo que el sistema no hace}
\begin{itemize}
    \item No reemplaza al intérprete humano de LSA: actúa como un ajuste
    razonable complementario, sobre un vocabulario acotado.
    \item No almacena datos clínicos ni el historial de conversaciones: el
    estado de la conversación es volátil (en memoria del navegador, vía
    BroadcastChannel).
    \item No envía datos a servidores externos fuera de la máquina local
    (el modelo de lenguaje y la síntesis de voz corren localmente).
    \item No reconoce LSA como lengua continua y completa: reconoce un
    conjunto acotado de clases de seña entrenadas.
\end{itemize}

\subsubsection{Exclusiones de alcance (TP2)}

La Declaración de Alcance oficial (TP2) excluye explícitamente del proyecto:
capacitación general en LSA al personal del sanatorio; desarrollo de
hardware a medida (solo ensamblaje de hardware comercial, ver E07 en el
EDT); integración con la Historia Clínica Electrónica (HCE) del sanatorio;
interpretación de modismos regionales complejos en esta primera versión;
mantenimiento técnico posterior a la entrega (15/11/2026); y
reestructuración de redes o cableado del sanatorio. Estas exclusiones
delimitan el alcance de este documento y de los entregables 01 a 09.

\subsubsection{Beneficios, objetivos y metas}
\begin{itemize}
    \item Reducir la barrera de comunicación en centros de salud sin
    intérprete de LSA disponible.
    \item Preservar la privacidad de los datos de salud mediante
    procesamiento local.
    \item Demostrar la viabilidad de un flujo extremo a extremo de visión
    por computadora, procesamiento de lenguaje natural y síntesis de voz
    en hardware de consumo.
\end{itemize}

\subsection{Definiciones, acrónimos y abreviaturas}

\begin{longtable}{p{4cm} p{10cm}}
    \caption{Definiciones, acrónimos y abreviaturas}
    \label{tab:definiciones} \\
    \toprule
    \textbf{Término} & \textbf{Definición} \\
    \midrule
    \endfirsthead
    \toprule
    \textbf{Término} & \textbf{Definición} \\
    \midrule
    \endhead
    LSA & Lengua de Señas Argentina --- lengua natural con sintaxis propia. \\
    LSTM & Long Short-Term Memory --- red neuronal recurrente para secuencias temporales, usada para clasificar señas. \\
    LLM & Large Language Model --- modelo de lenguaje grande (Qwen2 1.5B, servido localmente por Ollama). \\
    TTS & Text-to-Speech --- síntesis de texto a voz (motor Piper, modelo ONNX). \\
    MediaPipe Holistic & Framework de Google para extracción de landmarks corporales en tiempo real; opción de captura por defecto de LSA Salud. \\
    MediaPipe Tasks & Variante de MediaPipe basada en landmarkers independientes (Pose, Hand, Face); opción de captura alternativa, activable con \texttt{?mp=tasks}. \\
    Landmark & Punto anatómico detectado con coordenadas (x, y, z, visibilidad). \\
    Frame & Conjunto de 75 landmarks (33 de pose, 21 por mano) capturados en un instante. \\
    Blendshape & Coeficiente de expresión facial (0 a 1) provisto por MediaPipe Tasks; LSA Salud captura 52 por frame en modo Tasks --- base de captura para una futura incorporación de gestos no manuales/faciales a la clasificación (ver riesgo \#16, Sección~\ref{sec:riesgos-diseno}). \\
    Idle & Estado de reposo del paciente (manos abajo, sin movimiento significativo) que participa en el cierre de una oración. \\
    Plantilla médica & Frase médica predefinida asociada a una combinación conocida de clases de seña. \\
    Modo Captura & Modo de la Estación del Paciente para grabar nuevas muestras de entrenamiento. \\
    ERS / SRS & Especificación de Requisitos de Software / Software Requirements Specification. \\
    CU & Caso de Uso. \\
    RF / RNF & Requisito Funcional / Requisito No Funcional. \\
    REST & Representational State Transfer --- estilo arquitectónico para servicios web. \\
    WebSocket (WS) & Protocolo de comunicación bidireccional full-duplex sobre TCP. \\
    ONNX & Open Neural Network Exchange --- formato abierto para modelos de IA. \\
    BroadcastChannel API & API del navegador para comunicación entre pestañas del mismo origen, sin pasar por la red. \\
    p50 / p95 & Percentil 50 (mediana) y percentil 95 de una distribución de latencias. \\
    \bottomrule
\end{longtable}

\subsection{Referencias}

\begin{enumerate}
    \item IEEE Std 830-1998 --- IEEE Recommended Practice for Software Requirements Specifications.
    \item ISO/IEC/IEEE 42010 / UML 2.5 --- Object Management Group, \textit{Unified Modeling Language, Version 2.5.1}.
    \item Larman, C. \textit{UML y Patrones: Introducción al análisis y diseño orientado a objetos y al proceso unificado}. Pearson, 3.\textsuperscript{a} edición.
    \item Ley 26.378 --- Aprobación de la Convención sobre los Derechos de las Personas con Discapacidad (Argentina, 2008).
    \item Ley 27.044 --- Otorga jerarquía constitucional a la Convención (Argentina, 2014).
    \item Ley 25.326 --- Protección de Datos Personales (Argentina, 2000).
    \item Google MediaPipe --- \url{https://mediapipe.dev}
    \item PyTorch --- \url{https://pytorch.org}
    \item Ollama --- \url{https://ollama.ai}
    \item Piper TTS --- \url{https://github.com/rhasspy/piper}
    \item \texttt{README.md} y código fuente del repositorio ``Nexora'' (nombre interno del repositorio de código; \texttt{app.py}, \texttt{train\_lstm.py}, \texttt{src/\allowbreak model/\allowbreak lstm.py}, \texttt{web/}).
    \item TP0 --- Traductor LSA Salud (documento de tema); TP1 --- Acta de Constitución del Proyecto; TP2 --- Declaración de Alcance del Proyecto; Estructura de Desglose de Trabajo (EDT); Matriz de Trazabilidad de Requisitos; Registro de Interesados; Registro de Riesgos (documentos oficiales de cátedra, UTN-FRT).
\end{enumerate}

\subsection{Visión general del documento}

El resto de este documento se organiza así:
\begin{itemize}
    \item \textbf{Sección 2 --- Descripción general:} contexto del
    producto, funciones generales, características de los usuarios,
    restricciones y suposiciones.
    \item \textbf{Sección 3 --- Requisitos específicos:} interfaces
    externas, requisitos funcionales, requisitos de rendimiento,
    restricciones de diseño y atributos del sistema.
    \item \textbf{Apéndices:} matriz de trazabilidad entre requisitos
    funcionales, casos de uso y requisitos no funcionales.
\end{itemize}

% ============================================================================
\section{Descripción general}

\subsection{Perspectiva del producto}

LSA Salud es un sistema \textbf{independiente} (\textit{standalone}) que
corre íntegramente en una PC local. No depende de servicios en la nube
para su funcionamiento (el modelo de lenguaje y la síntesis de voz corren
localmente); si el proceso de Ollama no está disponible, la traducción cae
en un modo degradado que igual entrega las palabras originales
(ver RF-01.5, extensión 3a.2a).

El sistema usa una \textbf{arquitectura de doble estación}: la Estación
del Paciente (cámara, predicción y chat) y la Estación de Recepción
(historial de chat y respuesta por texto o voz), sincronizadas mediante
BroadcastChannel. Como se detalla en la Sección~2.5, estas dos estaciones
son, a la vez, el diseño de las dos caras de la estación fija de
guardia/triaje exigida por los requisitos oficiales RNF-04 y RNF-07.

\begin{table}[H]
    \centering
    \caption{Stack tecnológico de LSA Salud}
    \label{tab:stack}
    \small
    \begin{tabularx}{\textwidth}{>{\raggedright\arraybackslash}p{4cm} >{\raggedright\arraybackslash}p{4.4cm} X}
        \toprule
        \textbf{Componente} & \textbf{Tecnología} & \textbf{Función} \\
        \midrule
        Visión por computadora & MediaPipe Holistic / MediaPipe Tasks & Extrae landmarks corporales (y blendshapes faciales en modo Tasks). \\
        Clasificación de señas & PyTorch (LSTM) & Clasifica ventanas de 45 frames en clases de seña. \\
        Reconstrucción gramatical & Plantillas médicas + Ollama (Qwen2 1.5B) & Convierte palabras sueltas en una oración natural. \\
        Síntesis de voz & Piper TTS (ONNX) & Verbaliza la oración con voz argentina. \\
        Reconocimiento de voz & Web Speech API & Permite al personal de recepción o enfermería dictar respuestas. \\
        Sincronización entre estaciones & BroadcastChannel API & Mensajes entre pestañas del mismo origen, sin red. \\
        Frontend & Next.js 14 + React + TypeScript & Interfaz web de las estaciones. \\
        Backend & FastAPI + Uvicorn & API REST y WebSocket para inferencia y traducción. \\
        \bottomrule
    \end{tabularx}
\end{table}

\subsection{Funciones del producto}

\begin{enumerate}
    \item \textbf{Traducción de LSA a texto:} captura de señas, clasificación por LSTM y reconstrucción de oraciones naturales (CU-01, CU-02).
    \item \textbf{Síntesis de voz:} conversión de la oración traducida a audio con voz argentina (CU-03).
    \item \textbf{Canal de comunicación bidireccional:} chat sincronizado entre las dos estaciones (CU-01, CU-04).
    \item \textbf{Respuesta por texto o por voz:} el personal de recepción o enfermería responde escribiendo o dictando (CU-04, CU-05).
    \item \textbf{Captura de datos de entrenamiento:} grabación de nuevas señas para ampliar el vocabulario (CU-06, CU-07).
    \item \textbf{Entrenamiento del modelo:} reentrenamiento del LSTM con los datos capturados (CU-08).
    \item \textbf{Diseñado, pendiente de implementación:} consentimiento de cámara (CU-09), reporte de traducción incorrecta (CU-10) y registro de auditoría (CU-11).
\end{enumerate}

\subsection{Características de los usuarios}

\begin{table}[H]
    \centering
    \caption{Tipos de usuario y perfiles}
    \label{tab:usuarios}
    \small
    \begin{tabularx}{\textwidth}{>{\raggedright\arraybackslash}p{3.6cm} >{\raggedright\arraybackslash}p{4.1cm} X}
        \toprule
        \textbf{Actor} & \textbf{Perfil} & \textbf{Interacción} \\
        \midrule
        Paciente (ACT-01) & Usuario de LSA, sin conocimiento técnico esperado. & Hace señas frente a la cámara; lee las respuestas del profesional en pantalla. \\
        Personal de Recepción o Enfermería (ACT-02) & Personal administrativo o de enfermería. & Lee las traducciones del paciente; responde por texto o por dictado de voz. \\
        Desarrollador de IA / Equipo de Proyecto (ACT-03) & Personal técnico del equipo del proyecto. & Usa el Modo Captura para grabar señas nuevas y ejecuta el reentrenamiento. \\
        \bottomrule
    \end{tabularx}
\end{table}

\subsection{Restricciones}

\begin{itemize}
    \item \textbf{Ejecución local obligatoria:} ningún dato de la conversación sale de la máquina; el único tráfico de red es hacia procesos locales (Ollama en \texttt{localhost:11434}).
    \item \textbf{Dependencia de Ollama:} el servicio debe estar corriendo localmente para la reconstrucción gramatical por LLM; si no está disponible, se degrada a texto sin reconstruir (RF-01.5, extensión 3a.2a).
    \item \textbf{Navegador:} Google Chrome es necesario para Web Speech API (dictado) y para un soporte confiable de MediaPipe.
    \item \textbf{Vocabulario limitado:} el modelo entrenado reconoce actualmente 5 clases (\texttt{cabeza}, \texttt{dolor}, \texttt{idle}, \texttt{respirar}, \texttt{yo}); es ampliable mediante el Modo Captura y el reentrenamiento.
    \item \textbf{Sin persistencia de conversaciones:} todo el historial de chat es volátil (memoria del navegador + BroadcastChannel).
    \item \textbf{Lenguajes:} Python 3.9+ en el backend; TypeScript/React en el frontend.
    \item \textbf{Hardware mínimo:} PC con cámara web; GPU opcional para acelerar la inferencia del LSTM (funciona sobre CPU).
\end{itemize}

\subsection{Suposiciones y dependencias}

\begin{itemize}
    \item Se asume que el centro de salud dispone de al menos una PC con dos pantallas (o dos dispositivos) y una cámara web: las dos estaciones sincronizadas (Paciente y Recepción) implementan las dos caras de la estación fija de guardia/triaje exigida por RNF-04/RNF-07 oficiales (ver \texttt{IDS.md}, Sección~4.2); su ensamblaje físico sobre hardware comercial es tarea del entregable E07 del EDT, fuera de alcance de este documento.
    \item Se asume que Ollama y el modelo \texttt{qwen2:1.5b} están instalados y descargados de antemano.
    \item Se asume que el modelo de voz Piper (\texttt{es\_AR-daniela-high.onnx}) está disponible en la carpeta \texttt{voices/}.
    \item Se asume que el paciente conoce LSA y puede realizar señas de forma clara frente a la cámara.
    \item Se asume conexión eléctrica estable; no se contempla operación a batería.
    \item Si cambia el sistema operativo (probado en Windows), puede requerirse revalidar la compatibilidad de MediaPipe y Piper.
\end{itemize}

\subsection{Riesgos relevantes al diseño}
\label{sec:riesgos-diseno}

El Registro de Riesgos oficial (TP5) identifica riesgos con impacto directo
en requisitos de este documento; se listan los de calificación
\textit{Inaceptable} más relevantes al software:

\begin{itemize}
    \item \textbf{Riesgo \#1} (falta de disponibilidad de intérpretes
    oficiales de LSA para validar el dataset clínico, 30\%): afecta
    RNF-06.1 (validación lingüística).
    \item \textbf{Riesgo \#4} (iluminación deficiente en boxes que impida
    la captura de gestos, 25\%): afecta RF-01.1 (captura de landmarks).
    \item \textbf{Riesgo \#5} (latencia de procesamiento de video superior
    al umbral de 1 s, 35\%): afecta RNF-01.1 (rendimiento).
    \item \textbf{Riesgo \#9} (dificultad para compilar suficientes
    muestras de vocabulario médico, 30\%): afecta RNF-05.1 (especialización
    de datos).
    \item \textbf{Riesgo \#16} (falla del modelo al capturar gestos
    faciales/no manuales vitales en la sintaxis de LSA, 30\%): LSA Salud
    ya captura 52 \textit{blendshapes} faciales por frame en modo
    MediaPipe Tasks (ver Sección~2.4, Definiciones) como base de datos
    para una futura incorporación de gestos no manuales al modelo de
    clasificación; hoy esos coeficientes se capturan y se persisten
    (RF-T-04) pero todavía no se usan como entrada del LSTM, por lo que la
    mitigación está en curso y no cerrada.
\end{itemize}

\subsection{Requisitos futuros}

\begin{itemize}
    \item \textit{Requisito futuro:} ampliación del vocabulario a decenas de señas médicas frecuentes.
    \item \textit{Requisito futuro:} soporte para variantes regionales de LSA.
    \item \textit{Requisito futuro:} interfaz de administración web para gestionar modelos y datasets (hoy el reentrenamiento se hace por línea de comandos, CU-08).
    \item \textit{Requisito futuro:} verificación automática de que la traducción del LLM no incorpore síntomas no mencionados (hoy solo se indica por prompt, ver CU-02).
    \item \textit{Requisito futuro:} soporte de navegadores adicionales a Google Chrome para el dictado de voz.
    \item \textit{Requisito futuro:} restringir el origen permitido en CORS (hoy \texttt{allow\_origins=["*"]} en \texttt{app.py}), pensado para uso en una red local controlada.
    \item \textit{Requisito futuro:} integración opcional con sistemas de Historia Clínica Electrónica.
\end{itemize}

% ============================================================================
\section{Requisitos específicos}

Esta sección detalla los requisitos con nivel suficiente para el diseño,
la implementación y la verificación del sistema. Cada requisito está
identificado de forma unívoca y es trazable al código fuente (ver
Apéndice~A).

\subsection{Interfaces externas}

\subsubsection{Interfaces de usuario}

\begin{table}[H]
    \centering
    \caption{Interfaces de usuario del sistema}
    \label{tab:interfaces-usuario}
    \begin{tabularx}{\textwidth}{l l X}
        \toprule
        \textbf{Interfaz} & \textbf{Ruta} & \textbf{Descripción} \\
        \midrule
        Landing page & \texttt{/} & Selector de estación (Paciente / Recepción). \\
        Pantalla de Consentimiento (P-05) & \texttt{/paciente} (previa) & Consentimiento explícito antes de activar la cámara. \textit{Diseñada, pendiente de implementación.} \\
        Estación del Paciente & \texttt{/paciente} & Cámara, predicción en tiempo real, chat y Modo Captura. \\
        Estación de Recepción & \texttt{/recepcion} & Historial de chat, respuesta por texto y por dictado de voz. \\
        \bottomrule
    \end{tabularx}
\end{table}

\subsubsection{Interfaces de hardware}

\begin{table}[H]
    \centering
    \caption{Interfaces de hardware}
    \label{tab:interfaces-hw}
    \begin{tabularx}{\textwidth}{l l X}
        \toprule
        \textbf{Hardware} & \textbf{API} & \textbf{Uso} \\
        \midrule
        Cámara web & \texttt{navigator.mediaDevices.getUserMedia()} & Captura de video para MediaPipe. \\
        Micrófono & Web Speech API (\texttt{SpeechRecognition}) & Dictado de voz del personal de recepción o enfermería. \\
        Parlantes & \texttt{HTMLAudioElement} & Reproducción de la voz sintetizada (Piper TTS). \\
        \bottomrule
    \end{tabularx}
\end{table}

\subsubsection{Interfaces de software}

\begin{table}[H]
    \centering
    \caption{Endpoints del backend FastAPI}
    \label{tab:endpoints}
    \begin{tabularx}{\textwidth}{l l X}
        \toprule
        \textbf{Método} & \textbf{Endpoint} & \textbf{Descripción} \\
        \midrule
        GET & \texttt{/} & Health check: estado, dispositivo de cómputo, cantidad de clases. \\
        GET & \texttt{/classes} & Lista de clases que reconoce el modelo cargado. \\
        GET & \texttt{/dataset\_stats?mode=} & Cantidad de muestras por etiqueta, para el modo legado o Tasks. \\
        POST & \texttt{/predict} & Predicción de una seña vía REST (frames en JSON). \\
        WS & \texttt{/ws/predict} & Predicción de señas vía WebSocket (protocolo v1 JSON o v2 binario). \\
        POST & \texttt{/capture} & Guarda una muestra de entrenamiento (frames + etiqueta $\rightarrow$ \texttt{.npy}). \\
        POST & \texttt{/process\_sentence} & Traduce palabras sueltas a oración natural y genera el audio. \\
        GET & \texttt{/debug/timings.jsonl} & Transmite el registro de latencias por etapa (solo si está habilitado por variable de entorno; herramienta de desarrollo). \\
        \midrule
        POST & \texttt{/consent} \textit{(diseñado)} & Registra el consentimiento de cámara del paciente (CU-09, RF-03.1). \\
        POST & \texttt{/report} \textit{(diseñado)} & Registra un reporte de traducción incorrecta o ambigua (CU-10, RF-06.1). \\
        POST & \texttt{/audit} \textit{(diseñado, interno)} & Registra un evento de auditoría anónimo (CU-11, RF-04.1); no es invocado por el frontend, sino por el propio backend al completar cada operación de sistema. \\
        \bottomrule
    \end{tabularx}
\end{table}

También hay una dependencia de software externa: el servicio Ollama
(API REST local en el puerto 11434) que sirve el modelo de lenguaje
\texttt{qwen2:1.5b}.

\subsubsection{Interfaces de comunicación}

\begin{itemize}
    \item \textbf{BroadcastChannel:} canal \texttt{nexora-chat} (nombre interno heredado del repositorio de código) para sincronizar mensajes entre pestañas del mismo origen. Eventos: \texttt{MESSAGE} (nuevo mensaje) y \texttt{CLEAR} (limpiar historial), con deduplicación por identificador de mensaje.
    \item \textbf{WebSocket:} conexión persistente entre el frontend y \texttt{/ws/predict}, con reconexión automática cada 2 segundos ante corte.
    \item \textbf{HTTP REST:} comunicación solicitud-respuesta para traducción de oraciones, captura de datos y consulta de clases/estadísticas.
\end{itemize}

\subsection{Requisitos funcionales}

Los seis requisitos funcionales oficiales (RF-01 a RF-06, matriz de
trazabilidad TP2) se desarrollan aquí como sub-requisitos técnicos
trazables, numerados \texttt{RF-<oficial>.<n>} (ver \texttt{IDS.md},
Sección~3, para el texto literal de cada RF oficial y la tabla de
trazabilidad completa). Los requisitos sin contraparte oficial --- decisión
de diseño del equipo para CU-06/07/08 --- usan el prefijo \texttt{RF-T-}.
Cada uno indica su prioridad (\textit{Esencial}: imprescindible para el
flujo principal; \textit{Condicional}: mejora el flujo pero el sistema es
utilizable sin él; \textit{Deseable}: conveniencia menor) y los casos de
uso relacionados (ver Entregable de Modelos de Casos de Uso). Al final de
la sección se detallan RF-03.1, RF-04.1 y RF-06.1: \textbf{diseñados,
pendientes de implementación}.

\subsubsection{RF-01.1: Captura de landmarks en tiempo real (RF-01 oficial)}
\begin{table}[H]
\centering
\begin{tabularx}{\textwidth}{l X}
\toprule
\textbf{ID} & RF-01.1 \\
\textbf{Prioridad} & Esencial \\
\textbf{Descripción} & El sistema debe extraer, en tiempo real, los landmarks corporales del paciente desde la cámara web. \\
\textbf{Entrada} & Flujo de video de la cámara ($\sim$30 fps). \\
\textbf{Proceso} & MediaPipe (Holistic por defecto, o Tasks con \texttt{mp=tasks}) procesa cada frame elegible y extrae landmarks de pose (33) y de ambas manos (21 cada una), normalizados respecto de la nariz. \\
\textbf{Salida} & Un \textit{frame} de 75 landmarks (300 valores) por cada muestreo. \\
\textbf{CU relacionados} & CU-01. \\
\bottomrule
\end{tabularx}
\end{table}

\subsubsection{RF-01.2: Envío de frames para predicción (RF-01 oficial)}
\begin{table}[H]
\centering
\begin{tabularx}{\textwidth}{l X}
\toprule
\textbf{ID} & RF-01.2 \\
\textbf{Prioridad} & Esencial \\
\textbf{Descripción} & El frontend debe acumular frames en un buffer y enviarlos periódicamente al backend para su clasificación. \\
\textbf{Entrada} & Buffer de frames (hasta 45, con al menos 20 acumulados antes del primer envío). \\
\textbf{Proceso} & Cada 8 frames nuevos, el buffer se serializa como JSON (protocolo v1) o binario (protocolo v2, con coordenadas redondeadas) y se envía por WebSocket. \\
\textbf{Salida} & Mensaje WebSocket entrante con la predicción (ver RF-01.3). \\
\textbf{CU relacionados} & CU-01. \\
\bottomrule
\end{tabularx}
\end{table}

\subsubsection{RF-01.3: Clasificación de la seña mediante LSTM (RF-01 oficial; realiza también RF-05 oficial)}
\begin{table}[H]
\centering
\begin{tabularx}{\textwidth}{l X}
\toprule
\textbf{ID} & RF-01.3 \\
\textbf{Prioridad} & Esencial \\
\textbf{Descripción} & El backend debe clasificar la ventana de frames recibida en una de las clases de seña conocidas, con un nivel de confianza. El umbral de confianza (ver ``Umbral'' más abajo) es, a la vez, el mecanismo que realiza el requisito oficial RF-05 (``detectar cuando una seña no pueda interpretarse con confianza y evitar mostrar una traducción incorrecta''). \\
\textbf{Entrada} & Tensor de dimensiones $(1, 45, 300)$, recortado o relleno con ceros según haga falta. \\
\textbf{Proceso} & Red LSTM (2 capas, 128 unidades ocultas) seguida de una capa totalmente conectada y \textit{softmax}; se calcula la clase de mayor probabilidad y el top-3. \\
\textbf{Salida} & Seña predicha, confianza (\%) y top-3 de predicciones. \\
\textbf{Umbral} & $\geq$70\% de confianza para que la palabra se registre en la secuencia (ver RF-01.4). \\
\textbf{CU relacionados} & CU-01, CU-02. \\
\bottomrule
\end{tabularx}
\end{table}

\subsubsection{RF-01.4: Construcción de la secuencia y cierre de oración (RF-01 oficial; realiza también RF-05 oficial)}
\begin{table}[H]
\centering
\begin{tabularx}{\textwidth}{l X}
\toprule
\textbf{ID} & RF-01.4 \\
\textbf{Prioridad} & Esencial \\
\textbf{Descripción} & El sistema debe acumular las palabras reconocidas en una secuencia y detectar cuándo el paciente terminó de expresar su idea, descartando las señas de baja confianza o de reposo (mecanismo que también realiza RF-05 oficial). \\
\textbf{Entrada} & Predicciones sucesivas del LSTM y magnitud de movimiento entre frames consecutivos. \\
\textbf{Proceso} & Se agrega una palabra a la secuencia cuando su confianza es $\geq$70\% y no es la clase de reposo (\texttt{idle}) ni domina el top-3 con más de 50\%. La secuencia se cierra cuando la magnitud media de movimiento de los últimos 6 frames es menor a 0.05 y pasaron al menos 800 ms desde la última predicción distinta de reposo. \\
\textbf{Salida} & Secuencia de palabras cerrada, lista para traducirse (dispara CU-02). \\
\textbf{CU relacionados} & CU-01. \\
\bottomrule
\end{tabularx}
\end{table}

\subsubsection{RF-01.5: Traducción de la secuencia a oración natural (RF-01 oficial)}
\begin{table}[H]
\centering
\begin{tabularx}{\textwidth}{l X}
\toprule
\textbf{ID} & RF-01.5 \\
\textbf{Prioridad} & Esencial \\
\textbf{Descripción} & El sistema debe convertir la secuencia de palabras sueltas en una oración en español natural, priorizando respuestas instantáneas. \\
\textbf{Entrada} & Cadena de palabras sueltas (por ejemplo, ``yo dolor cabeza''). \\
\textbf{Proceso} & Se normaliza la secuencia y se busca, en orden: coincidencia exacta en el diccionario de plantillas médicas; luego en la caché de traducciones recientes (hasta 200 entradas); si no hay coincidencia, se delega en el modelo de lenguaje local (tiempo de espera de 8 segundos), con instrucciones explícitas de no inventar síntomas y de respetar la persona gramatical. \\
\textbf{Salida} & Oración natural (por ejemplo, ``Me duele la cabeza.''), y el origen de la traducción (plantilla, caché o modelo de lenguaje). \\
\textbf{CU relacionados} & CU-02. \\
\bottomrule
\end{tabularx}
\end{table}

\subsubsection{RF-01.6: Síntesis de voz (RF-01 oficial)}
\begin{table}[H]
\centering
\begin{tabularx}{\textwidth}{l X}
\toprule
\textbf{ID} & RF-01.6 \\
\textbf{Prioridad} & Esencial \\
\textbf{Descripción} & La oración traducida debe sintetizarse como audio, con voz en español argentino. \\
\textbf{Entrada} & Oración en español (texto). \\
\textbf{Proceso} & Si el audio ya fue generado antes para la misma oración, se reutiliza desde la caché (hasta 50 entradas); en caso contrario, se sintetiza con Piper TTS y el resultado se codifica en base64. \\
\textbf{Salida} & Audio WAV codificado en base64, reproducido en el navegador tras una ventana de exhibición del texto (1000 ms por defecto). \\
\textbf{CU relacionados} & CU-03. \\
\bottomrule
\end{tabularx}
\end{table}

\subsubsection{RF-01.7: Sincronización de la conversación entre estaciones (RF-01 oficial; soporta también RF-02 oficial)}
\begin{table}[H]
\centering
\begin{tabularx}{\textwidth}{l X}
\toprule
\textbf{ID} & RF-01.7 \\
\textbf{Prioridad} & Esencial \\
\textbf{Descripción} & Los mensajes del paciente y del personal de recepción o enfermería deben sincronizarse instantáneamente entre las dos estaciones. Es el canal que entrega el mensaje del profesional en un formato accesible para el paciente sordo (RF-02 oficial). \\
\textbf{Entrada} & Mensaje nuevo (del paciente, generado por CU-02/CU-03, o del personal, generado por CU-04). \\
\textbf{Proceso} & El mensaje se agrega al historial local y se publica por BroadcastChannel; cada estación deduplica por identificador antes de mostrarlo. \\
\textbf{Salida} & Historial de conversación actualizado en ambas estaciones. \\
\textbf{CU relacionados} & CU-01, CU-04. \\
\bottomrule
\end{tabularx}
\end{table}

\subsubsection{RF-02.1: Respuesta del personal por texto (RF-02 oficial)}
\begin{table}[H]
\centering
\begin{tabularx}{\textwidth}{l X}
\toprule
\textbf{ID} & RF-02.1 \\
\textbf{Prioridad} & Esencial \\
\textbf{Descripción} & El personal de recepción o enfermería debe poder escribir una respuesta y enviarla a la conversación compartida. \\
\textbf{Entrada} & Texto ingresado por el personal. \\
\textbf{Proceso} & Se valida que el texto no esté vacío antes de habilitar el envío. \\
\textbf{Salida} & Mensaje agregado a la conversación (ver RF-01.7). \\
\textbf{CU relacionados} & CU-04. \\
\bottomrule
\end{tabularx}
\end{table}

\subsubsection{RF-02.2: Dictado por voz del personal (RF-02 oficial)}
\begin{table}[H]
\centering
\begin{tabularx}{\textwidth}{l X}
\toprule
\textbf{ID} & RF-02.2 \\
\textbf{Prioridad} & Condicional \\
\textbf{Descripción} & El personal de recepción o enfermería debe poder dictar su respuesta por voz en lugar de escribirla. \\
\textbf{Entrada} & Audio del micrófono. \\
\textbf{Proceso} & Se usa la Web Speech API del navegador, configurada en español argentino, con resultados parciales y finales; el texto dictado se acumula sobre el texto ya escrito. \\
\textbf{Salida} & Texto transcripto, listo para enviarse (RF-02.1). \\
\textbf{Restricción} & Requiere un navegador compatible (Google Chrome). \\
\textbf{CU relacionados} & CU-05. \\
\bottomrule
\end{tabularx}
\end{table}

\subsubsection{RF-T-01: Selección o creación de etiqueta de seña (sin RF oficial; decisión de diseño del equipo)}
\begin{table}[H]
\centering
\begin{tabularx}{\textwidth}{l X}
\toprule
\textbf{ID} & RF-T-01 \\
\textbf{Prioridad} & Condicional \\
\textbf{Descripción} & Antes de grabar una muestra, el desarrollador debe poder elegir una etiqueta existente o crear una nueva. \\
\textbf{Entrada} & Nombre de etiqueta (existente o nueva). \\
\textbf{Proceso} & Si la etiqueta es nueva, se valida contra un patrón (letras minúsculas, dígitos y guion bajo, de 1 a 32 caracteres) antes de aceptarla; la carpeta de datos se crea recién al guardar la primera muestra. \\
\textbf{Salida} & Etiqueta activa para la siguiente grabación. \\
\textbf{CU relacionados} & CU-07. \\
\bottomrule
\end{tabularx}
\end{table}

\subsubsection{RF-T-02: Consulta de clases y estadísticas del dataset (sin RF oficial; decisión de diseño del equipo)}
\begin{table}[H]
\centering
\begin{tabularx}{\textwidth}{l X}
\toprule
\textbf{ID} & RF-T-02 \\
\textbf{Prioridad} & Condicional \\
\textbf{Descripción} & El sistema debe exponer las clases que reconoce el modelo y la cantidad de muestras grabadas por etiqueta. \\
\textbf{Entrada} & Modo activo (legado o Tasks). \\
\textbf{Proceso} & El backend lista las clases del modelo cargado y cuenta los archivos \texttt{.npy} por carpeta de etiqueta, en el árbol de datos correspondiente al modo. \\
\textbf{Salida} & Lista de clases y mapa etiqueta $\rightarrow$ cantidad de muestras. \\
\textbf{CU relacionados} & CU-07. \\
\bottomrule
\end{tabularx}
\end{table}

\subsubsection{RF-T-03: Grabación de muestras de entrenamiento (sin RF oficial; decisión de diseño del equipo)}
\begin{table}[H]
\centering
\begin{tabularx}{\textwidth}{l X}
\toprule
\textbf{ID} & RF-T-03 \\
\textbf{Prioridad} & Condicional \\
\textbf{Descripción} & El sistema debe permitir grabar una o varias muestras de una seña, con cuenta regresiva y repeticiones configurables. \\
\textbf{Entrada} & Etiqueta activa, cantidad de repeticiones. \\
\textbf{Proceso} & Por cada repetición: cuenta regresiva de 3 segundos, grabación de 45 frames consecutivos, envío al backend. \\
\textbf{Salida} & Confirmación de guardado y conteo de muestras actualizado por etiqueta. \\
\textbf{CU relacionados} & CU-06. \\
\bottomrule
\end{tabularx}
\end{table}

\subsubsection{RF-T-04: Persistencia de muestras de entrenamiento (sin RF oficial; decisión de diseño del equipo)}
\begin{table}[H]
\centering
\begin{tabularx}{\textwidth}{l X}
\toprule
\textbf{ID} & RF-T-04 \\
\textbf{Prioridad} & Condicional \\
\textbf{Descripción} & Cada muestra grabada debe guardarse en disco, separando el árbol de captura legado del basado en MediaPipe Tasks. \\
\textbf{Entrada} & Frames grabados (y, en modo Tasks, blendshapes faciales). \\
\textbf{Proceso} & Se valida que los frames sean numéricamente finitos y, en modo Tasks, que haya exactamente 52 blendshapes por frame; se guarda como archivo \texttt{.npy} bajo \texttt{dataset/custom/<etiqueta>/} (legado) o \texttt{dataset/custom\_tasks/<etiqueta>/} más \texttt{dataset/custom\_blendshapes/<etiqueta>/} (Tasks). \\
\textbf{Salida} & Archivo(s) \texttt{.npy} persistido(s) con nombre de marca de tiempo. \\
\textbf{CU relacionados} & CU-06. \\
\bottomrule
\end{tabularx}
\end{table}

\subsubsection{RF-T-05: Reentrenamiento del modelo LSTM (sin RF oficial; decisión de diseño del equipo)}
\begin{table}[H]
\centering
\begin{tabularx}{\textwidth}{l X}
\toprule
\textbf{ID} & RF-T-05 \\
\textbf{Prioridad} & Condicional \\
\textbf{Descripción} & El sistema debe permitir reentrenar el modelo LSTM con las muestras capturadas, fuera de línea. \\
\textbf{Entrada} & Conjunto de archivos \texttt{.npy} agrupados por etiqueta en \texttt{dataset/custom}. \\
\textbf{Proceso} & Descubrimiento de clases, separación 80/20 entrenamiento/validación (semilla fija), aumento de datos en entrenamiento (recorte aleatorio, ruido gaussiano, escalado), 50 rondas de entrenamiento, conservación del mejor modelo por exactitud de validación. \\
\textbf{Salida} & Archivo de pesos del modelo y archivo de clases actualizados. \\
\textbf{CU relacionados} & CU-08. \\
\bottomrule
\end{tabularx}
\end{table}

\subsubsection{RF-T-06: Selección de estación de trabajo (sin RF oficial; decisión de diseño del equipo)}
\begin{table}[H]
\centering
\begin{tabularx}{\textwidth}{l X}
\toprule
\textbf{ID} & RF-T-06 \\
\textbf{Prioridad} & Deseable \\
\textbf{Descripción} & Desde la página de inicio, el usuario debe poder abrir la estación que corresponda a su rol. \\
\textbf{Entrada} & Selección del usuario. \\
\textbf{Proceso} & Navegación a \texttt{/paciente} o \texttt{/recepcion}. \\
\textbf{Salida} & Estación correspondiente abierta. \\
\textbf{CU relacionados} & Precondición de CU-01 y de CU-04. \\
\bottomrule
\end{tabularx}
\end{table}

\subsubsection{RF-03.1: Consentimiento explícito de cámara (RF-03 oficial) --- diseñado, pendiente de implementación}
\begin{table}[H]
\centering
\begin{tabularx}{\textwidth}{l X}
\toprule
\textbf{ID} & RF-03.1 \\
\textbf{Estado} & \textbf{Diseñado --- pendiente de implementación.} \\
\textbf{Prioridad} & Esencial (bloqueante: sin consentimiento, no hay CU-01). \\
\textbf{Descripción} & Antes de activar la cámara, el sistema debe desplegar una pantalla de consentimiento explícito (P-05) y bloquear el acceso a la cámara hasta que el paciente lo otorgue. \\
\textbf{Entrada} & Decisión del paciente (acepta / no acepta). \\
\textbf{Proceso} & El sistema muestra el texto explicativo de P-05 antes de invocar \texttt{getUserMedia()}; si el paciente acepta, registra el consentimiento (sin datos identificatorios) y recién entonces solicita el permiso de cámara al navegador. \\
\textbf{Salida} & Consentimiento registrado y cámara activada, o cámara inactiva con alternativa ofrecida. \\
\textbf{CU relacionados} & CU-09 (precondición de CU-01). \\
\bottomrule
\end{tabularx}
\end{table}

\subsubsection{RF-04.1: Registro de eventos de auditoría (RF-04 oficial) --- diseñado, pendiente de implementación}
\begin{table}[H]
\centering
\begin{tabularx}{\textwidth}{l X}
\toprule
\textbf{ID} & RF-04.1 \\
\textbf{Estado} & \textbf{Diseñado --- pendiente de implementación.} \\
\textbf{Prioridad} & Condicional (no bloquea el flujo de comunicación; es una subfunción de auditoría). \\
\textbf{Descripción} & El sistema debe registrar, de forma automática y anónima, un evento de auditoría por cada operación de sistema relevante (predicción, traducción, síntesis, respuesta, captura), sin datos identificatorios ni imágenes o video, para habilitar métricas operativas de uso (E08, Informe de Impacto Operativo). \\
\textbf{Entrada} & Finalización de una operación de sistema (SO-01 a SO-05, ver Entregable 08). \\
\textbf{Proceso} & El backend arma un evento con tipo de operación, marca temporal y resultado, y lo persiste en el registro de auditoría basado en archivos (ver Entregable 09). \\
\textbf{Salida} & Evento de auditoría anónimo persistido. \\
\textbf{CU relacionados} & CU-11 (incluido por CU-01 y CU-04). \\
\bottomrule
\end{tabularx}
\end{table}

\subsubsection{RF-06.1: Reporte de traducción incorrecta o ambigua (RF-06 oficial) --- diseñado, pendiente de implementación}
\begin{table}[H]
\centering
\begin{tabularx}{\textwidth}{l X}
\toprule
\textbf{ID} & RF-06.1 \\
\textbf{Estado} & \textbf{Diseñado --- pendiente de implementación.} \\
\textbf{Prioridad} & Condicional. \\
\textbf{Descripción} & El paciente o el personal de recepción o enfermería deben poder marcar un mensaje de la conversación como traducción incorrecta o ambigua, indicando un motivo, para registrar el caso y priorizar el reentrenamiento (CU-08). \\
\textbf{Entrada} & Mensaje señalado y motivo del reporte. \\
\textbf{Proceso} & El sistema asocia el reporte al mensaje señalado y lo persiste, sin datos identificatorios del paciente. \\
\textbf{Salida} & Reporte registrado y confirmación mostrada al actor. \\
\textbf{CU relacionados} & CU-10 (extiende a CU-01 y CU-04). \\
\bottomrule
\end{tabularx}
\end{table}

\subsection{Requisitos de rendimiento (RNF-01.1, sub-requisito de RNF-01 oficial)}

Los siguientes valores son \textbf{objetivos de diseño}, expresados en
p50/p95 como exige la metodología del proyecto, y trazables al requisito
oficial RNF-01 (``latencia <1 segundo'', CRIT-01). El repositorio incluye un
mecanismo de instrumentación por etapa (variable de entorno
\texttt{LATENCY\_V2\_TIMING}, endpoint \texttt{/debug/timings.jsonl}) capaz
de medirlos, pero al momento de este documento \textbf{no hay una
medición formal registrada}; por eso la columna ``Estado'' distingue
entre objetivo y medición, y el criterio de aceptación oficial CRIT-01 no
puede darse todavía por verificado.

\begin{table}[H]
    \centering
    \caption{Requisitos de rendimiento (RNF-01.1)}
    \label{tab:rendimiento}
    \small
    \begin{tabularx}{\textwidth}{>{\raggedright\arraybackslash}p{5.2cm} >{\raggedright\arraybackslash}p{3.6cm} X}
        \toprule
        \textbf{Métrica} & \textbf{Valor objetivo} & \textbf{Estado} \\
        \midrule
        Latencia extremo a extremo (captura $\rightarrow$ audio listo), ruta con plantilla & p50 $\leq$ 600 ms & Objetivo de diseño \\
        Latencia extremo a extremo, caso general & p95 $\leq$ 1 s & Objetivo de diseño \\
        Inferencia LSTM por ventana (\texttt{lstm\_ms}) con MediaPipe Tasks vs.\ Holistic & sin regresión de p95 mayor al 10\% & Objetivo de diseño (gate manual documentado en \texttt{openspec}) \\
        Captura de frames & $\sim$30 fps (según cámara) & Dependiente de hardware \\
        Usuarios simultáneos & 1 par paciente--personal por instancia del backend & Restricción de diseño actual \\
        \bottomrule
    \end{tabularx}
\end{table}

\textbf{RNF-02.1 (sub-requisito de RNF-02 oficial, precisión):} la
exactitud de validación del modelo LSTM (\texttt{ModeloDeReconocimiento.exactitudDeValidación},
ver Entregable 01) se calcula durante \texttt{train\_lstm.py} sobre el 20\%
de validación, pero no existe todavía un harness formal que la reporte
como criterio de aceptación automatizado frente al 90\% exigido por
RNF-02/CRIT-02 oficiales. \textbf{Estado: objetivo de diseño, sin medición
formal registrada}; se recomienda incorporarlo al entregable de gestión E05
(Informe de Pruebas Controladas).

\subsection{Restricciones de diseño}

\begin{itemize}
    \item \textbf{Privacidad por diseño:} todo el procesamiento sensible ocurre en la máquina local; el único proceso externo invocado es Ollama, también local.
    \item \textbf{Backend de un solo archivo:} \texttt{app.py} concentra los endpoints REST y WebSocket, pensado para simplicidad de despliegue en un entorno de consultorio.
    \item \textbf{Modelos pre-entrenados como dependencia:} los pesos del LSTM, el modelo de voz Piper y el modelo de lenguaje de Ollama deben estar descargados de antemano; el sistema no los descarga en tiempo de ejecución.
    \item \textbf{Protocolos por tipo de operación:} WebSocket para predicción de baja latencia, REST para operaciones puntuales (traducción, captura, consulta de clases), BroadcastChannel para sincronización entre estaciones.
    \item \textbf{Sin base de datos:} las muestras de entrenamiento se guardan como archivos \texttt{.npy} en el sistema de archivos; la conversación es volátil y no se persiste.
\end{itemize}

\subsection{Atributos del sistema}

Los requisitos de rendimiento se documentan en la sección 3.3 (RNF-01.1).
Esta sección cubre los restantes atributos de calidad: los primeros tres
(RNF-03.1, RNF-04.1/RNF-07.1, RNF-05.1/RNF-06.1) son sub-requisitos
trazables a la matriz oficial; los últimos tres (RNF-T-01 a RNF-T-03) son
requisitos técnicos sin contraparte oficial.

\subsubsection{Seguridad y privacidad (RNF-03.1, sub-requisito de RNF-03 oficial)}
\begin{itemize}
    \item \textbf{Privacidad de datos de salud:} todo el procesamiento y almacenamiento ocurre en la máquina local; no hay persistencia de conversaciones (estado en memoria más BroadcastChannel).
    \item \textbf{Autenticación:} no implementada en el prototipo actual (\textit{Requisito futuro}).
    \item \textbf{Configuración de CORS:} el backend acepta actualmente cualquier origen (\texttt{allow\_origins=\allowbreak["*"]} en \texttt{app.py}); apropiado para pruebas en una red local controlada, pero debe restringirse antes de un despliegue de producción (\textit{Requisito futuro}, ver 2.7).
    \item \textbf{Acceso a cámara y micrófono:} requiere autorización explícita del navegador, otorgada por el usuario, y --- cuando RF-03.1 esté implementado --- consentimiento explícito previo del paciente (CU-09).
    \item \textbf{Eventos de auditoría (RF-04.1) y reportes de traducción (RF-06.1):} diseñados sin datos identificatorios ni imágenes o video, consistentes con esta restricción.
\end{itemize}

\subsubsection{Compatibilidad de hardware y estación fija (RNF-04.1 / RNF-07.1, sub-requisitos de RNF-04 y RNF-07 oficiales)}
\begin{itemize}
    \item Las dos estaciones sincronizadas (Paciente y Recepción), sobre el mismo navegador y canal BroadcastChannel, con doble monitor, implementan las dos caras de la estación fija de guardia/triaje que exigen RNF-04 y RNF-07 oficiales (cámara, pantalla doble cara, micrófono, altavoces, componentes comerciales estándar).
    \item \textbf{Estado: parcial.} El diseño de software está completo; el ensamblaje físico de las dos pantallas y la cámara sobre un mostrador o box existente, con hardware comercial, es tarea del entregable E07 ``Estación Desplegada en Guardia'' (EDT), fuera del alcance de software de este documento.
\end{itemize}

\subsubsection{Especialización y validación del dataset (RNF-05.1 / RNF-06.1, sub-requisitos de RNF-05 y RNF-06 oficiales)}
\begin{itemize}
    \item El Modo Captura y los árboles \texttt{dataset/custom}, \texttt{dataset/custom\_tasks}, \texttt{dataset/custom\_blendshapes} permiten ampliar el vocabulario a señas médicas; la cobertura completa del dominio clínico (síntomas, triaje, dolor) no está verificada al día de este documento (RNF-05.1, \textbf{estado: parcial}).
    \item La validación lingüística del dataset por intérpretes oficiales de LSA es un proceso organizacional (EDT, paquete 3.2.5), no una función de software (RNF-06.1, \textbf{estado: pendiente}); el riesgo \#1 del Registro de Riesgos aplica directamente (ver Sección~2.6).
\end{itemize}

\subsubsection{RNF-T-01: Fiabilidad (sin RNF oficial)}
\begin{itemize}
    \item El WebSocket de predicción se reconecta automáticamente cada 2 segundos ante un cierre de conexión.
    \item Si el modelo de lenguaje local falla o no responde a tiempo (8 s), se usa el texto original de la secuencia sin interrumpir el flujo.
    \item Si el motor Piper TTS no está cargado, el mensaje se entrega sin audio.
    \item Un indicador de solicitud en curso evita enviar frames al WebSocket mientras una predicción anterior sigue pendiente.
\end{itemize}

\subsubsection{RNF-T-02: Disponibilidad (sin RNF oficial)}
\begin{itemize}
    \item El sistema no depende de conectividad a internet para el flujo crítico (captura, clasificación LSTM, síntesis de voz): todos esos componentes corren localmente.
    \item La disponibilidad de la traducción por modelo de lenguaje depende de que el proceso local de Ollama esté activo; su ausencia degrada la traducción pero no interrumpe el sistema (ver RNF-T-01).
\end{itemize}

\subsubsection{RNF-T-03: Mantenibilidad (sin RNF oficial)}
\begin{itemize}
    \item Separación clara entre frontend (Next.js) y backend (FastAPI).
    \item Responsabilidades divididas en hooks de React independientes: captura de cámara, WebSocket de predicción, construcción de oraciones, canal de chat y reconocimiento de voz.
    \item El modelo LSTM está desacoplado del backend en un módulo propio (\texttt{src/\allowbreak model/\allowbreak lstm.py}).
    \item Las plantillas médicas están centralizadas en una estructura de datos del backend, fácil de extender.
\end{itemize}

\subsubsection{RNF-T-04: Portabilidad (sin RNF oficial)}
\begin{itemize}
    \item Probado en Windows; el stack (Next.js, FastAPI/Uvicorn, PyTorch, ONNX) es multiplataforma en principio.
    \item Requiere Node.js 18 o superior y Python 3.9 o superior.
\end{itemize}

\subsubsection{RNF-T-05: Usabilidad y accesibilidad (sin RNF oficial)}
\begin{itemize}
    \item Conversación fluida, sin botones de confirmación manual: la ventana de exhibición del texto antes de reproducir el audio actúa como confirmación automática (excepción deliberada: el consentimiento de cámara de P-05/CU-09 y el reporte de traducción de CU-10 sí requieren una acción explícita, ver Entregable 04).
    \item Interfaz en español, con indicadores visuales de estado (``EN VIVO'', ``GRABANDO'', cuenta regresiva, barras de confianza en tiempo real).
    \item Arquitectura de doble estación pensada para separar visualmente el rol del paciente y el del profesional de salud.
\end{itemize}

\subsection{Otros requisitos}

\subsubsection{Escalabilidad del vocabulario}
\begin{itemize}
    \item El Modo Captura permite ampliar el vocabulario de forma incremental, sin requerir conocimientos técnicos avanzados (CU-06, CU-07).
    \item El reentrenamiento es offline mediante \texttt{train\_lstm.py} (CU-08).
\end{itemize}

\subsubsection{Compatibilidad}
\begin{itemize}
    \item Navegador: Google Chrome (requerido para Web Speech API y recomendado para MediaPipe).
    \item Sistema operativo: Windows (probado); Linux/macOS teóricamente compatibles, sin validación registrada.
\end{itemize}

% ============================================================================
\section{Apéndices}

\subsection{Apéndice A: Matriz de trazabilidad RF -- CU -- RNF}

\small
\begin{longtable}{>{\raggedright\arraybackslash}p{3.4cm} >{\raggedright\arraybackslash}p{5.6cm} >{\raggedright\arraybackslash}p{5.6cm}}
    \caption{Matriz de trazabilidad Requisitos Funcionales -- Casos de Uso -- Requisitos No Funcionales}
    \label{tab:trazabilidad} \\
    \toprule
    \textbf{Requisito funcional} & \textbf{Casos de uso relacionados} & \textbf{Requisitos no funcionales relacionados} \\
    \midrule
    \endfirsthead
    \toprule
    \textbf{Requisito funcional} & \textbf{Casos de uso relacionados} & \textbf{Requisitos no funcionales relacionados} \\
    \midrule
    \endhead
    RF-01.1 & CU-01 & RNF-T-01 \\
    RF-01.2 & CU-01 & RNF-T-01, RNF-01.1 \\
    RF-01.3 & CU-01, CU-02 & RNF-01.1 \\
    RF-01.4 & CU-01 & RNF-T-05, RNF-01.1 \\
    RF-01.5 & CU-02 & RNF-T-01, RNF-01.1 \\
    RF-01.6 & CU-03 & RNF-T-01, RNF-01.1 \\
    RF-01.7 & CU-01, CU-04 & RNF-T-05 \\
    RF-02.1 & CU-04 & --- \\
    RF-02.2 & CU-05 & RNF-T-04 (compatibilidad de navegador) \\
    RF-T-01 & CU-07 & --- \\
    RF-T-02 & CU-07 & --- \\
    RF-T-03 & CU-06 & --- \\
    RF-T-04 & CU-06 & RNF-03.1 \\
    RF-T-05 & CU-08 & --- \\
    RF-T-06 & Precondición de CU-01, CU-04 & --- \\
    \midrule
    RF-03.1 \textit{(diseñado)} & CU-09 & RNF-03.1 \\
    RF-04.1 \textit{(diseñado)} & CU-11 & RNF-03.1 \\
    RF-06.1 \textit{(diseñado)} & CU-10 & RNF-03.1 \\
    \bottomrule
\end{longtable}

\small
\begin{longtable}{>{\raggedright\arraybackslash}p{3.4cm} >{\raggedright\arraybackslash}p{10.2cm}}
    \caption{Matriz de trazabilidad RNF oficiales -- sub-requisitos}
    \label{tab:trazabilidad-rnf} \\
    \toprule
    \textbf{RNF oficial} & \textbf{Sub-requisitos técnicos} \\
    \midrule
    \endfirsthead
    \toprule
    \textbf{RNF oficial} & \textbf{Sub-requisitos técnicos} \\
    \midrule
    \endhead
    RNF-01 (latencia) & RNF-01.1 \\
    RNF-02 (precisión) & RNF-02.1 \\
    RNF-03 (seguridad/privacidad) & RNF-03.1 \\
    RNF-04 (hardware) & RNF-04.1 (compartido con RNF-07.1) \\
    RNF-05 (dataset especializado) & RNF-05.1 \\
    RNF-06 (validación lingüística) & RNF-06.1 \\
    RNF-07 (estación fija) & RNF-07.1 (compartido con RNF-04.1) \\
    \bottomrule
\end{longtable}

\subsection{Apéndice B: Formatos de datos relevantes}

\subsubsection{Formato de un frame (entrada al pipeline de reconocimiento)}
\begin{verbatim}
Frame = [
  [x_0, y_0, z_0, vis_0],     // landmark 0 (pose)
  ...                          // 33 landmarks de pose
  [x_33, y_33, z_33, vis_33],  // landmark 33 (mano izquierda)
  ...                          // 21 landmarks de mano izquierda
  [x_54, y_54, z_54, vis_54],  // landmark 54 (mano derecha)
  ...                          // 21 landmarks de mano derecha
]
// Total: 75 landmarks x 4 valores = 300 flotantes por frame
\end{verbatim}

\subsubsection{Formato de respuesta de predicción (WebSocket)}
\begin{verbatim}
{
  "sign": "dolor",
  "confidence": 88.5,
  "top3": [
    {"sign": "dolor", "confidence": 88.5},
    {"sign": "cabeza", "confidence": 8.2},
    {"sign": "idle", "confidence": 3.3}
  ]
}
\end{verbatim}

\subsubsection{Formato de respuesta de traducción (REST)}
\begin{verbatim}
{
  "original": "yo dolor cabeza",
  "natural_sentence": "Me duele la cabeza.",
  "audio_base64": "UklGRi4AAABXQVZFZm10IBAA...",
  "source": "template",
  "latency_ms": 3
}
\end{verbatim}
